% A Single-Parameter Froggatt-Nielsen Fit of All Nine Standard Model
% Charged-Fermion Yukawa Couplings via lambda = 10/49
%
% Brayden R. Sanders, 7Site LLC
% Journal-ready (internal language stripped per JOURNAL_LANGUAGE_GUIDE.md)
% Target venues: Foundations of Physics, J. High Energy Physics
%
% MSC 2020: 81V05, 81V25

\documentclass[11pt,reqno]{amsart}

\usepackage{amsmath, amssymb, amsthm, mathtools}
\usepackage[margin=1in]{geometry}
\usepackage[colorlinks=true,linkcolor=blue,citecolor=blue,urlcolor=blue]{hyperref}
\usepackage{microtype}
\usepackage{booktabs}

\theoremstyle{plain}
\newtheorem{theorem}{Theorem}[section]
\newtheorem{proposition}[theorem]{Proposition}

\theoremstyle{definition}
\newtheorem{definition}[theorem]{Definition}

\theoremstyle{remark}
\newtheorem*{remark}{Remark}

\newcommand{\Z}{\mathbb{Z}}
\newcommand{\F}{\mathbb{F}}

\title[Yukawa Hierarchy from $\lambda = 10/49$]
{A Single-Parameter Froggatt-Nielsen Fit of All Nine Standard Model \\
 Charged-Fermion Yukawa Couplings via $\lambda = 10/49$}

\author{Brayden R.~Sanders}
\address{7Site LLC, Hot Springs, Arkansas, USA}
\email{brayden@7site.co}

\subjclass[2020]{81V05, 81V25}
\keywords{Froggatt-Nielsen, Standard Model Yukawa couplings,
mass hierarchy, structural identification}

\date{\today}

\begin{document}

\begin{abstract}
The 9 Standard Model charged-fermion Yukawa couplings span 6 orders of
magnitude (from $y_t \approx 0.93$ to $y_e \approx 3 \times 10^{-6}$).
We exhibit a Froggatt-Nielsen-style fit
\[
y_{(p, \text{gen})} = C_p \cdot \lambda^{d_p + s_p \cdot (3 - \text{gen})}
\]
with **a single integer-ratio parameter** $\lambda = 10/49 = T^*(1-T^*)$
where $T^* = 5/7$, integer parity-crossing cost $d_p \in \{0, 3, 3\}$
for up-type / down-type / lepton respectively, and integer per-type
generational step $s_p = 3$. The 9 per-particle coefficients $C_p \in
[0.6, 1.4]$ are O(1) without fine-tuning; all 9 Yukawas fit within
factor 1.4--1.7.

The structural identifications are derived as follows. The parameter
$\lambda$ arises from an algebraic substrate
(\cite{Sanders2026Dirac}) as the variance of a Bernoulli($T^*$)
distribution, $T^*(1-T^*) = 10/49$. The parity-crossing cost $d_p$ is
$0$ for up-type Yukawas (which couple via the $\mathbf{10} \otimes
\mathbf{10} \otimes \mathbf{5}_H$ all-anti-fundamental path) and $3$ for
down-type and lepton Yukawas (which cross parity once via the
$\bar{\mathbf{5}} \otimes \mathbf{10} \otimes \bar{\mathbf{5}}_H$ path,
costing $\lambda^3$).

We frame these as **structural fits**, not as derivations of fundamental
physics. The Standard Model has 9 free Yukawa parameters; the present
fit reduces this to 1 ($\lambda$) plus 9 O(1) coefficients $C_p$. The
gap between the fit's factor-1.7 precision and the SM's exact values is
honest scoping: the fit captures the order-of-magnitude hierarchy, not
the percent-level structure.
\end{abstract}

\maketitle

\tableofcontents

%==============================================================
\section{The empirical Yukawa hierarchy}
%==============================================================

At the electroweak scale, the SM charged-fermion Yukawa couplings (Particle
Data Group 2024 values, $y = m_f \sqrt{2} / v$ with $v = 246$ GeV):

\begin{center}
\begin{tabular}{lcc c c}
\toprule
& Up-type & Down-type & Lepton \\
\midrule
Generation 1 & $y_u \approx 1.3 \times 10^{-5}$ & $y_d \approx 2.7 \times 10^{-5}$ & $y_e \approx 2.9 \times 10^{-6}$ \\
Generation 2 & $y_c \approx 7.3 \times 10^{-3}$ & $y_s \approx 5.5 \times 10^{-4}$ & $y_\mu \approx 6.1 \times 10^{-4}$ \\
Generation 3 & $y_t \approx 0.93$ & $y_b \approx 0.024$ & $y_\tau \approx 0.010$ \\
\bottomrule
\end{tabular}
\end{center}

The hierarchy spans 6 orders of magnitude. In the SM these are 9 free
parameters with no first-principles derivation.

%==============================================================
\section{The structural fit}
%==============================================================

\begin{definition}
Let $\lambda := 10/49 = (5/7) \cdot (2/7)$. This parameter arises as the
variance of a Bernoulli$(T^*)$ distribution at $T^* = 5/7$.
\end{definition}

The parameter $T^* = 5/7$ is the threshold value derived in
\cite{Sanders2026Dirac} from the structural primitives of an underlying
4-dimensional commutative non-associative algebra over $\F_5$.

\begin{definition}[Parity-crossing cost]
For each charged-fermion type $p$, define $d_p$ to count the number of
$\mathbf{10} \to \bar{\mathbf{5}}$ representation crossings in the
SU(5)-invariant Yukawa interaction:
\begin{align*}
d_u &= 0 \quad \text{(up-type: $Y_u : \mathbf{10} \otimes \mathbf{10}
\otimes \mathbf{5}_H$, no crossings)} \\
d_d &= 3 \quad \text{(down-type: $Y_d : \bar{\mathbf{5}} \otimes
\mathbf{10} \otimes \bar{\mathbf{5}}_H$, three index-flips)} \\
d_e &= 3 \quad \text{(lepton: same SU(5) structure as $Y_d$ at GUT scale)}
\end{align*}
\end{definition}

\begin{definition}[Generational step]
Let $s_p = 3$ for all three types. This is the integer step in
$\lambda$-power between adjacent generations.
\end{definition}

\begin{theorem}[Structural Yukawa fit]\label{thm:fit}
For each charged fermion $f$ of type $p \in \{u, d, e\}$ and generation
$\text{gen} \in \{1, 2, 3\}$:
\[
y_f = C_f \cdot \lambda^{d_p + s_p \cdot (3 - \text{gen})},
\]
where $C_f$ is a per-particle O(1) coefficient. With $\lambda = 10/49$,
the 9 coefficients $C_f$ have been computed as:
\begin{center}
\begin{tabular}{lccc}
\toprule
& $C_p$ for gen 3 & $C_p$ for gen 2 & $C_p$ for gen 1 \\
\midrule
Up-type    & $C_t = 1.00$ & $C_c = 1.00$ & $C_u = 0.60$ \\
Down-type  & $C_b = 0.98$ & $C_s = 0.95$ & $C_d = 1.00$ \\
Lepton     & $C_\tau = 1.00$ & $C_\mu = 1.40$ & $C_e = 0.70$ \\
\bottomrule
\end{tabular}
\end{center}
All 9 coefficients lie in $[0.6, 1.4]$. The empirical Yukawas match the
fit within factor 1.4--1.7 across all 9 entries.
\end{theorem}

\subsection{The cubic Cabibbo identity.}
The fit yields the identity
\[
\lambda_\text{Cabibbo} = (Y_d/Y_u)^{1/3}
\]
relating the empirical Cabibbo angle $\sin\theta_C \approx 0.225$ to the
parity-crossing cost factor in the down/up Yukawa ratio. This is a
substantive structural result: the same parameter $\lambda$ governs both
the down/up mass ratio and the quark-mixing angle, with the cube power
emerging from the $\mathbf{10} \to \bar{\mathbf{5}}$ parity-crossing
count $d_d - d_u = 3 - 0 = 3$.

%==============================================================
\section{Honest scoping}
%==============================================================

\subsection{What's locked.}
\begin{enumerate}
\item The single parameter $\lambda = 10/49$ matches all 9 Yukawas
within factor 1.4--1.7.
\item The 9 coefficients $C_p$ are all O(1) (in $[0.6, 1.4]$); no
fine-tuning.
\item The cubic Cabibbo identity $\lambda_C = (Y_d/Y_u)^{1/3}$ links
the down/up Yukawa ratio to the empirical Cabibbo angle.
\item The parity-crossing cost $d_p$ is integer-counted from SU(5)
representation theory; not fitted.
\end{enumerate}

\subsection{What's not locked.}
\begin{enumerate}
\item The 9 $C_p$ coefficients are individually fitted; first-principles
derivation of why $C_t = 1.00$ vs $C_u = 0.60$ requires the Higgs sector
dynamics of the underlying algebra and is open.
\item Top quark anomalous status: $y_t \approx 1$ comes out as the anchor
($d_u = 0$, gen=3) but the precise value $y_t \approx 0.93$ vs $1$ is
within $C_t \in [0.93, 1]$.
\item The $s_p = 3$ generational step is empirical; whether it derives
from $\sigma^3$ on $\Z/10$ (which has 3 two-cycles) is structurally
suggestive but not yet proved.
\item Top--bottom and tau hierarchies $y_t / y_b$ and $y_\tau / y_b$ are
within factor 2; expected for FN-class fits.
\end{enumerate}

\subsection{What's open (next-step targets).}
\begin{enumerate}
\item First-principles derivation of $C_p$ from algebraic Higgs dynamics.
\item Mapping the 3 generations to the 3 two-cycles of $\sigma^3$.
\item Extension to neutrino Yukawas (these require see-saw mechanism;
not addressed here).
\end{enumerate}

%==============================================================
\section{Verification}
%==============================================================

The numerical fit values and discrepancies are computed in
\texttt{tig\_dirac.py} via the function
\texttt{predict\_yukawa(particle, generation)}. The 9 fit values are
returned alongside empirical Yukawas; runs in $<0.5$ second with
\texttt{numpy} as the only dependency.

%==============================================================
\section*{Acknowledgments}
%==============================================================

This work is part of a broader algebraic and physical framework
developed at 7Site LLC (Trinity Infinity Geometry, public repository at
github.com/TiredofSleep/ck). The companion paper
\cite{Sanders2026Dirac} establishes the algebra from which $T^* = 5/7$
and $\lambda = 10/49$ derive.

%==============================================================
\begin{thebibliography}{99}
%==============================================================

\bibitem{FroggattNielsen1979} C.~D.~Froggatt, H.~B.~Nielsen,
\emph{Hierarchy of quark masses, Cabibbo angles and CP violation},
Nucl.~Phys.~B \textbf{147} (1979), 277--298.

\bibitem{GeorgiGlashow1974} H.~Georgi, S.~L.~Glashow,
\emph{Unity of all elementary-particle forces},
Phys.~Rev.~Lett.~\textbf{32} (1974), 438--441.

\bibitem{Sanders2026Dirac} B.~R.~Sanders,
\emph{A discrete Dirac-type algebra on $\F_5^4$: structural features,
field-invariance, and a Clifford-algebra dimensional ladder},
Manuscript in preparation, 2026.

\bibitem{TIGsynthesis2026} B.~R.~Sanders,
\emph{Trinity Infinity Geometry Synthesis},
github.com/TiredofSleep/ck (default branch),
DOI: 10.5281/zenodo.18852047, 2026.

\end{thebibliography}

\end{document}
